% =============================================================================
\chapter{Pipeline Workspace: x121 — Adult Content}
\label{ch:pipeline-workspace}
% =============================================================================

Selecting a pipeline from the landing page enters the pipeline workspace: a
re-scoped sidebar with an added \emph{Pipeline Admin} group, and every Content
/ Production / Review / Tools route filtered to the chosen pipeline.

We walk \texttt{x121 — Adult Content} end to end. The drill-down example uses
the \textbf{SDG} project, and within SDG the \textbf{Amouranth} avatar.

\screenshot{09-ws-enter.png}{Pipeline card click — routing into the \texttt{x121} workspace.}
\screenshot{09-ws-sidebar.png}{Sidebar re-scoped to the \texttt{x121} workspace (note the Pipeline header and Pipeline Admin group).}

\begin{designnote}
The re-scoping is driven by \texttt{PipelineProvider} which reads
\texttt{\$pipelineCode} from the URL. All subsequent queries include the
pipeline ID, so even pages whose component is shared with the global nav
(e.g. Media, Scenes) show pipeline-filtered data.
\end{designnote}

% --------------------------------------------------------------------------
\section{Workspace Dashboard}
\label{sec:ws-dashboard}
% Route: /pipelines/x121/dashboard
% Component: app/pages/PipelineDashboardPage.tsx

Pipeline-scoped landing page. Same widget framework as the global dashboard, but
every query is pre-filtered to \texttt{x121}, and the header surfaces pipeline
seed-slot summary.

\screenshot{09-ws-dashboard-hero.png}{Pipeline dashboard for \texttt{x121}.}

\begin{itemize}
  \pagelement{Pipeline header}{Code, name, description, seed-slot summary.}
  \pagelement{Widget grid}{Pipeline-scoped Active Tasks and Project Progress widgets.}
  \pagelement{Quick Navigation}{Grid of card links to prominent pages within the workspace.}
\end{itemize}

% --------------------------------------------------------------------------
\section{All Projects → SDG}
\label{sec:ws-projects}
% Route: /pipelines/x121/projects
% Component: features/projects/ProjectListPage.tsx

Projects are the top-level container of avatars. The Project Hub list plus the
project detail page make the project the frontend's explicit working context.


\screenshot{09-ws-projects-list.png}{Pipeline-scoped project list.}

\subsection{SDG — project detail}
% Route: /pipelines/x121/projects/1
% Component: features/projects/ProjectDetailPage.tsx

Project detail is the central hub for avatar work. Its primary view is the
project's avatars; additional tabs cover scene enablement, batch production
state, delivery, and project settings.

\screenshot{09-ws-sdg-detail-hero.png}{SDG detail — default (Overview) tab.}

\begin{itemize}
  \pagelement{Breadcrumb}{Projects $\rightarrow$ SDG.}
  \pagelement{Tab strip}{Overview, Avatars, Production, Delivery, Settings.}
  \pagelement{Review assignments link}{Opens the assignment dashboard.}
\end{itemize}

\subsubsection{Overview tab}
Stat tickers plus the deliverables grid. Per-avatar readiness (metadata, images,
scenes, speech) is surfaced as coloured indicators on each card so the reader
can see at-a-glance which avatars still need work without drilling in.
\screenshot{09-ws-sdg-overview-tab.png}{SDG / Overview — stat tickers and deliverables grid.}
The AvatarDeliverablesGrid has two inner tabs:
\screenshot{09-ws-sdg-overview-readiness.png}{Overview / Readiness inner tab.}
\screenshot{09-ws-sdg-overview-matrix.png}{Overview / Matrix inner tab.}

\subsubsection{Avatars tab}
All avatars in the project. Avatars are pipeline-scoped: the same avatar name
imported into \texttt{x121} and \texttt{y122} is two completely independent
records with no cross-pipeline linking.
\screenshot{09-ws-sdg-avatars-tab.png}{SDG / Avatars tab — all avatars in the project.}

\subsubsection{Production tab}
Project-level view over the batch orchestrator's scene matrix — what is queued,
running and blocked across this project's avatars.
\screenshot{09-ws-sdg-production-tab.png}{SDG / Production tab.}

\subsubsection{Delivery tab}
Delivery readiness per avatar, keyed off the project's \texttt{blocking\_deliverables}
setting (non-blocking sections start deselected).
\screenshot{09-ws-sdg-delivery-tab.png}{SDG / Delivery tab.}

\subsubsection{Settings tab}
Project-level overrides: scene enablement, naming-rule overrides, delivery
output-profile default, metadata template inheritance.
\screenshot{09-ws-sdg-settings-tab.png}{SDG / Settings tab.}

% --------------------------------------------------------------------------
\section{SDG → Amouranth (avatar detail)}
\label{sec:ws-amouranth}
% Route: /pipelines/x121/projects/1/avatars/5
% Component: features/avatars/AvatarDetailPage.tsx

Amouranth is an example avatar inside SDG. The avatar detail page is the full
"workstation": the entire avatar workflow — image upload, seed assignment,
scene generation, review, metadata, speech and delivery — is achievable from
this single page without navigating away. The detail page has ten tabs.


\screenshot{09-amouranth-hero.png}{Amouranth detail — default (Overview) tab.}

\subsection{Overview}
Aggregated snapshot of the avatar's configuration and readiness — what is set,
what is missing, and a direct link into the fix-it flow for each gap.
\screenshot{09-amouranth-overview-tab.png}{Amouranth / Overview tab.}

\subsection{Images}
Per-track image grid. "Image" here means a media variant: the underlying
tables cover images, video and audio uniformly after the image-to-media
rename.
\screenshot{09-amouranth-images-tab.png}{Amouranth / Images tab — per-track image grid.}

\subsection{Seeds}
Seed slot cards, driven by the pipeline's \texttt{seed\_slots} definition (not
hardcoded clothed / topless). Each slot is a Workflow Media Slot — bound to a
\texttt{LoadImage} / \texttt{LoadVideo} / \texttt{LoadAudio} node in the workflow
JSON — and auto-detection picks the best matching approved variant by track
affinity.
\screenshot{09-amouranth-seeds-tab.png}{Amouranth / Seeds tab — seed slot cards.}

\subsection{Scenes}
Every scene video version for this avatar, with clip gallery, per-clip QA
state, and resume-from-last-good-clip for partial regenerations that avoid
re-rendering approved work.
\screenshot{09-amouranth-scenes-tab.png}{Amouranth / Scenes tab.}

\subsection{Derived}
Externally-processed clips cut from this avatar's approved scenes, with parent-
version lineage preserved. Same review tooling as generated clips; import is
available via browser upload or server-side directory scan.
\screenshot{09-amouranth-derived-tab.png}{Amouranth / Derived clips tab.}

\subsection{Metadata}
Dual-metadata editor backed by pipeline-scoped templates. Each pipeline defines
its own required metadata fields (x121 wants body type / ethnicity;
y122 wants speaking style / topic expertise).
\screenshot{09-amouranth-metadata-tab.png}{Amouranth / Metadata tab.}

\subsection{Speech}
Avatar dialogue strings organised by speech type and language — the text feed
for ElevenLabs TTS. Speech types are pipeline-scoped, so each pipeline has its
own set (x121 has "Flirty" / "Whisper"; y122 has "Introduction" / "Q\&A Response").

\screenshot{09-amouranth-speech-tab.png}{Amouranth / Speech tab.}

\subsection{Deliverables}
Four sections — Metadata, Images, Scene Videos, Speech — gated by the project's
\texttt{blocking\_deliverables} setting. Non-blocking sections start deselected
("ignored") so it's clear which deliverables the project actually requires.
\screenshot{09-amouranth-deliverables-tab.png}{Amouranth / Deliverables tab.}

\subsection{Review}
Avatar-level review gate: a holistic sign-off that says "this avatar, across all
its scenes and deliverables, is complete". Sits above the per-clip review and
can be auto-allocated with round-robin load balancing.

\screenshot{09-amouranth-review-tab.png}{Amouranth / Review tab.}

\subsection{Settings}
Per-avatar overrides on top of project / pipeline / catalogue defaults:
model settings, voice ID, scene enablement, and metadata overrides.
\screenshot{09-amouranth-settings-tab.png}{Amouranth / Settings tab.}

% --------------------------------------------------------------------------
\section{Pipeline-scoped Content}
\label{sec:ws-content}

Each of these mirrors the global Content pages in Chapter~\ref{ch:content}
but scoped to \texttt{x121}. Only the filter defaults and visible data
differ; the UI is otherwise identical. Internal links within the workspace
preserve the pipeline prefix so sidebar context isn't lost.

\screenshot{09-ws-avatars.png}{\texttt{x121} → Avatars browse.}
\screenshot{09-ws-library.png}{\texttt{x121} → Library.}
\screenshot{09-ws-media.png}{\texttt{x121} → Media.}
\screenshot{09-ws-scenes.png}{\texttt{x121} → Scenes.}
\screenshot{09-ws-derived-clips.png}{\texttt{x121} → Derived Clips.}
\screenshot{09-ws-catalogue.png}{\texttt{x121} → Catalogue (first tab).}

% --------------------------------------------------------------------------
\section{Pipeline-scoped Queue \& Review}
\label{sec:ws-queue-review}

Queue, Annotations and Reviews filtered to jobs, clips and assignments that
belong to \texttt{x121}.

\screenshot{09-ws-queue.png}{\texttt{x121} → Queue.}
\screenshot{09-ws-annotations.png}{\texttt{x121} → Annotations.}
\screenshot{09-ws-reviews.png}{\texttt{x121} → Reviews.}

% --------------------------------------------------------------------------
\section{Pipeline-scoped Tools}
\label{sec:ws-tools}

Workflows filtered to the pipeline. A scene type can only reference tracks and
workflows that belong to its own pipeline, so cross-pipeline leakage is caught
at save time.

\screenshot{09-ws-workflows.png}{\texttt{x121} → Workflows.}

% --------------------------------------------------------------------------
\section{Pipeline Admin}
\label{sec:ws-admin}

Two admin pages are scoped to the pipeline workspace sidebar. They target the
pipeline-level configuration that admins routinely edit without leaving the
workspace.

\subsection{Naming Rules (pipeline-scoped)}
% Route: /pipelines/x121/naming
% Component: features/naming-rules/NamingRulesPage.tsx (PipelineNamingRulesEditor)

Edits the pipeline's import matching rules and naming-rule JSONB — the tokens
and patterns that keep the naming engine aware of this pipeline's specifics
(e.g. x121's \texttt{clothed.png} / \texttt{topless.png} seed convention).


\screenshot{09-ws-naming-hero.png}{Pipeline-scoped naming editor — different UX to the admin view.}

\begin{itemize}
  \pagelement{Editor pane}{Edit the pipeline's \texttt{naming\_rules} JSONB directly.}
  \pagelement{Validation}{Inline errors for invalid tokens.}
  \pagelement{Save}{Captured, never submitted.}
\end{itemize}

\subsection{Output Profiles (pipeline-scoped)}
% Route: /pipelines/x121/output-profiles

Pipeline-level output-profile default. Sits in the inheritance chain between
the platform default and any project-specific override.

\screenshot{09-ws-output-profiles-hero.png}{Pipeline-scoped output profiles.}
